Cálculos numéricos
A.1.  Suma de números
A.2.  Resta de números
A.3.  Multiplicación de números
A.4.  División de números
A.5.  Resolución de problemas
Conceptos básicos de números negativos
A.6.  Ubicación en la recta numérica
A.7.  Comparación de negativos
A.8.  Oraciones con números negativos
A.9.  Determina el signo con sumas y restas
A.10. Determina el signo con multiplicaciones y divisiones
Operaciones con números negativos
A.11. Comparación de negativos
A.12. Suma y resta con negativos
A.13. Multiplicación y división con negativos
A.14. Potencias con números negativos
A.15. Jerarquía de operaciones
Fracciones MCM y MCD
A.16. Comparación de fracciones
A.17. Fracciones equivalentes
A.18. M.C.D y M.C.M
A.19. Simplificación de fracciones
A.20. Resolución de problemas
Operaciones con fracciones 
A.21. Suma y resta con denominadores iguales
A.22. Suma y resta denominadores diferentes
A.23. Multiplicación de fracciones
A.24. División de fracciones
A.25. Resolución de problemas
Leyes de los exponentes
A.26. Suma de exponentes
A.27. Resta de exponentes
A.28. Multiplicación de exponentes
A.29. División de exponentes
A.30. Exponentes negativos
Sucesiones aritméticas
A.31. Completando la sucesión
A.32. Diferencia de una sucesión
A.33. Término general
A.34. Término enésimo
A.35. Suma de una sucesión aritmética

Probabilidad y estadística
A.36. Mediana y moda
A.37. Promedio
A.38. Desviación media
A.39. Eventos mutuamente excluyentes
A.30. Eventos dependientes e independientes
Figuras y cuerpos geométricos
A.41. Perímetro
A.42. Área
A.43. Área lateral y total
A.44. Volumen
A.45. Resolución de problemas
Plano cartesiano y recta
A.46. Ubicación en el plano cartesiano
A.47. Cuadrantes en el plano cartesiano
A.48. Pendiente y ordenada
A.49. Pendiente dados dos puntos
A.50. Ecuación de una recta
Ecuación lineal
A.51. Lenguaje algebraico
A.52. Ecuaciones lineales 1
A.53. Ecuaciones lineales 2
A.54. Ecuaciones lineales con fracciones
A.55. Resolución de problemas
Sistemas de ecuaciones
A.56. Método de eliminación
A.57. Método de sustitución e igualación
A.58. Sistema de ecuaciones 3x3
A.59. Sistema de ecuaciones con fracciones
A.60. Resolución de problemas

Sucesiones cuadráticas y geométricas
A.61. Sucesión cuadrática
A.62. Completando la sucesión cuadrática
A.63. Término general
A.64. Sucesión geométrica
A.65. Razón de una sucesión geométrica
Factorización
A.66. Término común
A.67. Diferencia de cuadrados
A.68. Trinomio cuadrado perfecto
A.69. Trinomios de la forma $x^2+bx+c$
A.70. Trinomios de la forma $ax^2+bx+c$
Productos notables
A.71. Binomios conjugados
A.72. Binomios con término común
A.73. Binomio al cuadrado
A.74. Binomios de la forma $(mx+a)(nx+b)$
A.75. Binomio al cubo
Ecuaciones cuadráticas
A.76. Clasificación de ecuaciones cuadráticas
A.77. Discriminante
A.78. Ecuaciones cuadráticas incompletas
A.79. Ecuaciones cuadráticas completas 1
A.80. Ecuaciones cuadráticas completas 2
Teorema de Pitágoras
A.81. Identificación de lados
A.82. Hallando la hipotenusa
A.83. Hallando el cateto
A.84. Áreas y perímetros
A.85. Resolución de problemas
Trigonometría
A.86. Identificando lados
A.87. Identificando funciones
A.88. Encontrando lados
A.89. Encontrando ángulos
A.90. Resolución de problemas